% Generated draft: requires mathematical and source review.
\documentclass[11pt]{article}
\usepackage[margin=25mm]{geometry}
\usepackage{fontspec}
\setmainfont{Latin Modern Roman}
\usepackage{amsmath,amssymb,mathtools,unicode-math}
\setmathfont{Latin Modern Math}
\usepackage{enumitem,needspace,xurl,hyperref}
\setlength{\parindent}{0pt}
\setlength{\parskip}{5pt}
\setlength{\emergencystretch}{3em}
\newcommand{\N}{\mathbb{N}}
\newcommand{\Z}{\mathbb{Z}}
\newcommand{\Q}{\mathbb{Q}}
\newcommand{\R}{\mathbb{R}}
\newcommand{\C}{\mathbb{C}}
\newcommand{\F}{\mathbb{F}}
\newcommand{\A}{\mathbb{A}}
\newcommand{\PP}{\mathbb{P}}
\newcommand{\E}{\mathbb{E}}
\newcommand{\eps}{\varepsilon}
\newcommand{\dd}{\,\mathrm{d}}
\begin{document}
\section*{UnsolvedMath ID: 30001139}
\textbf{Source problem code:} catalogue entry
\subsection{Definitions and mathematical statement}
The supplied record groups three distinct questions under permanent UnsolvedMath ID \(30001139\), source code \texttt{catalogue entry}. They are retained separately below. Several hypotheses needed for fully specified mathematical problems are absent from the supplied record.

\begin{enumerate}
\item \textbf{Origami hyperbolic paraboloid.} Let \(S\) be a square sheet carrying the described crease pattern of concentric squares and diagonals. A no-stretch realization would be a continuous map \(f:S\to\R^3\) whose restriction to each face preserves intrinsic lengths, with each prescribed crease mapped to a straight segment and each face realized as a developable surface isometric to its original planar face. Does the intended folded hyperbolic-paraboloid model admit such a realization, or does realizing that model require stretching?

The record does not specify the number or positions of the concentric squares, the complete crease pattern, mountain and valley assignments, boundary conditions, permitted self-contact, or the conditions distinguishing the intended folded configuration from the flat sheet. It also does not specify whether ``hyperbolic paraboloid'' names an exact target surface or the familiar origami shape. Consequently, the existence question remains conditional on a precise specification of that model; the conditions above alone do not characterize it.

\item \textbf{A second triangle-area relation.} Let \(S=[0,1]^2\) be partitioned into \(T\) triangles \(\Delta_1,\ldots,\Delta_T\), where \(T\) is a positive integer. For an ordered triangle with vertices \(p_i,q_i,r_i\in\R^2\), define its signed area by
\[
 A_i=\frac12\det(q_i-p_i,r_i-p_i).
\]
With triangle orientations compatible with the positive orientation of the square, the stated area relation is
\[
 \sum_{i=1}^{T}A_i=1.
\]
What is the ``second relation'' among \(A_1,\ldots,A_T\), besides this sum relation?

The supplied record does not identify a fixed triangulation, restrictions on vertices, the intended sign convention beyond the displayed normalization, or the class of relations sought. It does not say whether the additional relation may depend on the partition's combinatorics, or what independence from the sum relation means. Thus it does not determine a specific additional identity or a fully quantified conjecture asserting one.

\item \textbf{Periodic inversive-distance packing rigidity.} For circles with centers \(c_u,c_v\in\R^2\) and positive radii \(r_u,r_v\), the usual unsigned inversive distance is
\[
 I(u,v)=\frac{\lVert c_u-c_v\rVert^2-r_u^2-r_v^2}{2r_ur_v}.
\]
A circle configuration is doubly periodic if two linearly independent translations preserve it, including its prescribed adjacency and edge labels. Consider the source's hexagonal-lattice circle packing with neighboring inversive distances prescribed by parameters \(a,b,c\). Is every such doubly periodic packing rigid? The source presents the following as an equivalent question: must every other packing with those prescribed distances have a common radius, that is,
\[
 \exists r>0\ \forall v,\qquad r_v=r?
\]

The record does not specify whether the adjacency graph is the triangular lattice of a hexagonal circle arrangement or the honeycomb graph, how \(a,b,c\) label its edges, or their admissible ranges. It also leaves unspecified the allowed circle intersections, the period lattice and its permitted variation, the comparison class of other packings, and whether rigidity is local or global and modulo which transformations. The inversive-distance formula above records a conventional meaning; the source's convention requires confirmation. The asserted equivalence between rigidity and the common-radius property is preserved as a source claim, not established here. These omissions prevent a unique precise rigidity conjecture.
\end{enumerate}
\subsection{Short English statement}
The record asks three separate questions: whether a specified origami saddle can be folded without stretching, what additional relation constrains triangle areas in a square partition, and whether certain doubly periodic circle packings are rigid. Essential specifications are missing in each formulation. The supplied open-status assessment and claimed packing equivalence require human review.
\subsection{Sources}
\begin{itemize}
\item Local UnsolvedMath registry, permanent ID \(30001139\), source code \texttt{catalogue entry}, title: \emph{Origami, Triangle-Area, and Circle-Packing Rigidity Problems}; collection: \emph{Oberwolfach Reports Open Problems}. The record gives proposed year \(2009\) and names no individual proposer.
\item Supplied report citation: \emph{Discrete Differential Geometry} (2009), Oberwolfach Reports, pp.~75--144, DOI \texttt{10.4171/owr/2009/02}. \url{https://doi.org/10.4171/owr/2009/02}. This is the supplied report-level reference; question-specific pages are not supplied.
\item Supplied literature source: Demaine et al., ``nonexistence of pleated folds.'' \url{https://arxiv.org/abs/0906.4747}. This is the supplied descriptive label; no more specific bibliographic attribution or theorem application is asserted here.
\item Supplied literature source: Rote, \emph{Open Problems in Discrete Differential Geometry}. \url{https://page.mi.fu-berlin.de/rote/Kram/OWR-DDG09-problems.pdf}.
\item Literature assessment dated exactly \texttt{2026-08-21}: the supplied curation labels the record open, says it combines an origami question with a separate inversive-distance packing rigidity question, and states that the latter is not resolved by the cited partial rigidity theorems. The background describes this as dated literature triage derived from supplied Oberwolfach Reports curation and requests independent verification of time-sensitive status and cited claims. This is supplied dated curation, not an independently verified literature assessment.
\item The registry labels the statement \texttt{corrected\_verified} and reports that all three extracted questions are identifiable and that report-preface prose was removed. These are supplied editorial assertions, not independent verification of the formulations or their mathematical status.
\end{itemize}
\end{document}
